\documentclass[11pt, letterpaper]{article}

\input{/home/hyperion/.config/LatexHub/preambles/base.tex}
\input{/home/hyperion/.config/LatexHub/preambles/diagrams.tex}
\input{/home/hyperion/.config/LatexHub/preambles/tikz.tex}
\input{/home/hyperion/.config/LatexHub/preambles/macros-cat-theory.tex}
\input{/home/hyperion/.config/LatexHub/preambles/macros-algebra.tex}
\input{/home/hyperion/.config/LatexHub/preambles/basic-theorems-section.tex}
\input{/home/hyperion/.config/LatexHub/preambles/refs-basic.tex}
\input{/home/hyperion/.config/LatexHub/preambles/homework-formatting.tex}
\usepackage{titlepageBU}
\title{Final Project Notes}
\begin{document}
\makereport
\section{Inner products or smt}
Given an inner product $\left\langle -,- \right\rangle : V\times V \to k$, we can identify it with a 2-cotensor $g\in \mathcal{T} ^2(V)$. Let $\left\{ v_i \right\} _{i=1} ^n\subseteq V$ a basis and $\left\{ v_i^* \right\} _{i=1} ^n \subseteq V^*$ the corresponding dual basis. Writing $g_{ij} =\left\langle v_i,v_j \right\rangle $, we have
\[
	\left\langle -,- \right\rangle =\sum_{i,j=1} ^n g_{ij} v^*_i\otimes v^*_j
.\]
We can see also that $\det g_{ij} \neq 0$.

\begin{proposition}\label{1.1}
	Let $V$ be an $n$-dimensional vector space with inner product $\left\langle -,- \right\rangle $. Then there is a basis $\left\{ v_i \right\} _{i=1} ^n$ with $\left\langle v_i,v_j \right\rangle =\delta _{ij} $; such a basis is called \emph{orthonormal}.
\end{proposition}
The proof is by Grahm-Schmidt. A nice corollary is that there is an isomorphism $f : \mathbb{R} ^n \to V$ such that for all $a,b\in\mathbb{R}^n $, the inner product in $\mathbb{R} ^n$, denoted by the dot product, can be given as
\[
	a\cdot b=\left\langle f(a),f(b) \right\rangle 
.\]
We write $g^{ij} =g_{ij} ^{-1} $. We apply this to manifolds.

\begin{definition}\label{1.2}
	Let $\pi : E \to M$ be a vector bundle. A \emph{Riemannian metric} on $E$ is a function $\left\langle -,- \right\rangle : M \to E$ which maps each $p\in M$ to a positive-definite inner product $\left\langle -,- \right\rangle _p : \pi ^{-1} (p) \to \mathbb{R} $, and is continuous in the sense that for any two continuous sections $s_1,s_2$ of $E$, the map $p\mapsto \left\langle s_1(p),s_2(p) \right\rangle_p$ is continuous $M\to \mathbb{R} $. If $M$ is a smooth manifold and $E$ is a smooth vector bundle
\end{definition}
\begin{theorem}\label{1.3}
	Let $\pi : E \to M$ be a smooth $k$-plane bundle over $M$. Then there is a smooth Riemannian metric on $E$.
\end{theorem}

When the bundle is the tangent bundle $TM$, a smooth Riemannian metric for $TM$ will simply be called a Riemannian metric for $M$. We can write this in coordinates as
\[
	\left\langle -,- \right\rangle =g_{ij} \mathrm{d} x^i\otimes \mathrm{d} x^j
.\]

\section{IDK}

Let $\gamma : [a,b] \to M$ be a smooth curve. This yields tangent vectors $\gamma '(t)\in T_{\gamma (t)} M$, and we can thus define the \emph{length} of the tangent vectors by
\[
	\left\| \frac{\mathrm{d}\gamma }{\mathrm{d}t}  \right\| =\sqrt{\left\langle \frac{\mathrm{d}\gamma }{\mathrm{d}t} ,\frac{\mathrm{d}\gamma }{\mathrm{d}t}  \right\rangle_{\gamma (t)} } 
.\]
We thus define the \emph{length of $\gamma $ from $a$ to $b$ as}
\[
	L_a^b(\gamma )=\int_{a}^{b} \left\| \frac{\mathrm{d}\gamma }{\mathrm{d}t}  \right\|  \,\mathrm{d} t=\int_{a}^{b} \left\| \gamma '(t) \right\|  \,\mathrm{d} t
.\]
If $\gamma $ is piecewise smooth, we define its length by summing over lengths of smooth components. We generally will simply write $L_a^b$ instead of $L$.

Classically, the norm $\left\| - \right\| $ was denoted by $\mathrm{d} s$. This notation is to be interpreted as $\left| \mathrm{d} s \right| =\gamma ^*(\left\| - \right\| )$.

We can define a metric on $M$ via $d(p,q)$ defined as the inf over lengths of smooth curve from $p$ to $q$. However, the shortest possible curve bewteen two points need not be a straight line. For example, if the manifold is $\mathbb{R} ^n$ with a hole at the origin, then the curve has to go around the hole. To find the true path, we employ methods from the calculus of variations.


For example, suppose there is a suitably differentiable function $F : \mathbb{R} ^3 \to \mathbb{R} $. We look for functions $f : [a,b] \to \mathbb{R} $ that minimize the quantity
\[
	J=\int _a^b F(t,f(t),f'(t)) \,\mathrm{d} t
.\]
To find minima and maxima for $J$, we have to look at curves in the set of functions. We proceed by taking a variation -- that is, a function $\alpha : (-\varepsilon ,\varepsilon )\times [a,b] \to \mathbb{R} $ such that $\alpha (0,t)=f(t)$. Thus $\alpha $ determines a family of functions in $t$ paramterized over $(-\varepsilon ,\varepsilon )$, and this family passes through $f$. We will frequently abuse notation and write $\alpha (u)=\alpha (u,-)$. If $\alpha $ also fixes endpoints in the sense that $\alpha(-,a)=f(a)$ and $\alpha (-,b)=f(b)$, then we say $\alpha $ is a variation of $f$ that fixes endpoints.

we now compute
\[
	\left.\frac{\mathrm{d}J(\alpha (u))}{\mathrm{d}u} \right_{u=0} =\left. \frac{\mathrm{d}}{\mathrm{d}u} \right|_{u=0} \int_{a}^{b} F(t,\alpha (u,t),\frac{\partial \alpha }{\partial t} (u,t)) \,\mathrm{d} t
.\]
\begin{align*}
	\left.\frac{\mathrm{d}J(\alpha (u))}{\mathrm{d}u} \right_{u=0} &=\left. \frac{\mathrm{d}}{\mathrm{d}u} \right|_{u=0} \int_{a}^{b} F(t,\alpha (u,t),\frac{\partial \alpha }{\partial t} (u,t)) \,\mathrm{d} t\\
								       &=\int_{a}^{b} \left. \frac{\mathrm{d}}{\mathrm{d}u} \right_{u=0} F(t,\alpha (u,t),\frac{\partial \alpha }{\partial t} (u,t)) \,\mathrm{d} t
.\end{align*}
?????
\[
	\Gamma _{ij} ^k=\frac{1}{2} \left( \frac{\partial g_{ik} }{\partial x^j} +\frac{\partial g_{jk} }{\partial x^i} -\frac{\partial g_{ij} }{\partial x^k}  \right) 
.\]
So
\[
	\frac{\mathrm{d}^2\gamma ^k}{\mathrm{d}t^2} +\sum_{i,j=1} ^n \Gamma _{ij} ^k \frac{\mathrm{d}\gamma ^i}{\mathrm{d}t} \frac{\mathrm{d}\gamma ^j}{\mathrm{d}t} =0
.\]

\end{document}
